\chapter{Fundamental group of a circle}

\section{The statement}

\begin{thm}{Theorem}\label{thm:pi1(S1)}
The fundamental group of the circle $\mathbb{S}^1$ is isomorphic to the additive group of integers $\mathbb{Z}$.
\end{thm}

Given a loop $f\:[0,1]\to \mathbb{S}^1$, denote by $\deg f$ the number of times it goes aroond $\mathbb{S}^1$ counteclockwise (each clockwise turn is counded with minus).
In the following sections we will define $\deg$ formally, prove that it depends only on the homotopy class of $f$ (in other words, if $f_0\sim f_1$, then $\deg f_0=\deg f_1$) and then show that $\deg\:\pi_1(\SS^1)\to\ZZ$ is an isomorphism.

Let us show how to apply this theorem before diving into the proof.

\begin{thm}{Exercise}\label{ex:S1}
Apply the theorem to show the following.

\begin{subthm}{ex:S1:R3>R^2}
The space $\RR^3$ is not homeomorphic to the plane $\RR^2$.
\end{subthm}

\begin{subthm}{ex:D-->S}
The circle $\SS^1=\partial \DD^2$ is not a retract of the disc $\DD^2$.
\end{subthm}

\begin{subthm}{ex:bry->bry}
Given $x\in\DD^2$ the complement $\DD^2\setminus\{x\}$ is simply connected if and only if $x\in \SS^1=\partial \DD^2$.
Conclude that any homeomorphism $\DD^2\to\DD^2$ maps $\partial \DD$ to itself.
\end{subthm}

\begin{subthm}{ex:S1:finite-complement}
The complement of any finite nonempty set $F\subset\RR^2$ is not contractible.
\end{subthm}

\end{thm}

Part \ref{SHORT.ex:D-->S} in the last exercise has the following application.

\begin{thm}{Brouwer fixed point theorem in the plane}
\label{cor:brouwer-D2}
Any continuous map $f\:\DD^2\to\DD^2$ has a fixed point.
\end{thm}

\parit{Proof.}
Arguing by contradiction, assume $x\neq f(x)$ for all $x\in\DD^2$.
Consider given $x\in\DD^2$ consider the half-line $H$ from $f(x)$ to $x$.
Note that $H$ intesects $\SS^1=\partial \DD^2$ at a sinle point; denote it by $h(x)$.
Note that $h$ is continuous and $h(x)=x$ if $x\in\partial \DD^2$;
that is, $h$ defines a retruction $\DD^2\to \SS^1$.
The later contradicts \ref{ex:D-->S}.
\qeds

\begin{wrapfigure}[6]{r}{50 mm}
\vskip-6mm
\centering
\includegraphics{mppics/pic-1405}
\end{wrapfigure}

\begin{thm}{Advanced exercise}\label{ex:KK}
Consider two sets $K_1$ and $K_2$ in the plane, shown in the picture.
Each one being a closed annulus with two attached line segments.
In $K_1$ one segment is attached from the inside and another from the outside;
in $K_2$ both segments are attached from the outside.

\begin{subthm}{ex:K!=K}
Show that $K_1 \not\simeq K_2$.
\end{subthm}

\begin{subthm}{ex:K=K}
Show that $K_1 \times[0,1]\simeq K_2 \times[0,1]$.
\end{subthm}

\end{thm}


\section{Liftings}

Let us identify $\RR^2$ with the complex plane $\CC$;
namely, we will encode a point $(x,y)\in\RR^2$ by the complex number $z=x+i\cdot y$; here $i$ denotes the imaginary unit.
In particular, $\SS^1=\set{z=x+i\cdot y\in \CC}{|z|^2=x^2+y^2=1}$.

Consider the map $e\:\RR\to \mathbb{S}^1\subset \CC$ defined by
\[e\:x\mapsto \exp(2\cdot\pi\cdot i\cdot x)=\cos(2\cdot\pi\cdot x)+i\cdot\sin(2\cdot\pi\cdot x)).\]
(The last equality is called Euler's idenity, it can be taken as a definition of $\exp(2\cdot\pi\cdot i\cdot x)$.)

Let $f\:\c{X}\to \SS^1$ be a continuous map.
A continuous map $\tilde f\:\c{X}\to \RR$ will be called \index{lifting}\emph{lifting} of $f$ if $f=e\circ \tilde f$;
in other words if the following diagram \index{commutative diagram}\emph{commutes};
\begin{figure}[h!]
\centering
\begin{tikzcd}
& \RR \arrow[d, "e"]\\
\c{X} \arrow[r, "f"']\arrow[ru, "\tilde f"] & \SS^1
\end{tikzcd}
\end{figure}
the latter means that following directed paths in the diagram with the same start and endpoints lead to the same result.
Even in this simple case, using a commutative diagram is helpful, and later on you simply can't do without them.

\begin{thm}{Proposition}\label{prop:lifting}
Let $f_0\:[0,1]\to\SS^1$ be a path and let $x_0\in \RR$ be a point such that $e(x_0)=f_0(0)$.
Then there is unidue lifting $\tilde f_0$ of $f$ such that $\tilde f_0(0)=x_0$.
Moreover, if $f_0\sim f_1$ then $\tilde f_0\sim \tilde f_1$.
\end{thm}

\begin{thm}{Lemma}\label{lem:covering}
Let $W=\SS^1\setminus \{u_0\}$ for some $u_0=e(x_0)\in\SS^1$.
Then $e^{-1}(W)$ is a disjoint union of open intervals $V_n=(x_0+n,x_0+n+1)$ for $n\in \ZZ$
and $e$ defines a homeomorphism $V_n\to U$ for each $n$.

In particular, given a connected space $\c{X}$, continuous map $f\:\c{X}\z\to W$, and a point $q\in \RR$ such that $e(q)=f(p)$ for some $p\in \c{X}$ there is a unique lifting $\tilde f\:\c{X}\to \RR$ of $f$ such that $\tilde f(p)=q$.
\end{thm}


\parit{Proof.}
Observe that $e_n=e|_{V_n}$ is homeomorphism $V_n\to W$.
Indeed, $e_n$ is a continuous bijection.
Furthermore, by \ref{obs:compact-to-hausdorff} $e|_{\bar V_n}$ is a closed map.
Therefore $e_n=e|_{V_n}$ is a closed continuous bijection $V_n\to W$ and hence a homeomorphism.

Note that the intervals $V_n$ are connected components of $e^{-1}W$.
In particular, $q\in V_n$ for some $n$.
Since $\c{X}$ is connected, $\tilde f(\c{X})\subset V_n$ for any lifting $\tilde f$ lifting of $f$ such that $e(q)=f(p)$.
Since $e_n=e|_{V_n}$ is a homeomorpism, we have $\tilde f=e_n^{-1}\circ f$, which proves existance and uiqueness.
\qeds


\parit{Proof of \ref{prop:lifting}.}
Consider two open subsets $W_1=\SS^1\setminus (1,0)$ and $W_2=\SS^1\setminus (-1,0)$.
Note that $\SS^1=W_1\cup W_2$.
Therefore, $[0,1]=f^{-1}(W_1)\z\cup f^{-1}(W_2)$.

Let $\eps>0$ be the Lebesgue number of this covering; see \ref{prop:lebesgue-number}.
Consider a particition $0=t_0<t_1<\dots<t_n=1$ of $[0,1]$ into equal intervals shorter than $\eps$.
Note that $f([t_i,t_{i+1}])\subset W_1$ or  $f([t_i,t_{i+1}])\subset W_2$ for each $i$.

By \ref{lem:covering}, there is unique lifting of $f|_{[t_0,t_1]}$ such that $\tilde f(0)=x_0$; let $x_1=\tilde f(t_1)$.
By the same argument, there is unique lifting of $f|_{[t_1,t_2]}$ such that $\tilde f(t_1)=x_1$.
Repeating this argument finitely many times produces a lift of $f$, which is unique by construction.
The continuity of $f$ follows from \ref{ex:u-continuous:closed}.

\begin{wrapfigure}{r}{30 mm}
\vskip-6mm
\centering
\includegraphics{mppics/pic-1410}
\end{wrapfigure}

The proof of the last statement is similar.
Choose a homotopy $F\:[0,1]\times [0,1]\to\SS^1$ from $f_0$ to $f_1$.
Applying \ref{prop:lebesgue-number} again, we can subdivide $[0,1]\times [0,1]$ into small squares by horizontal and vertical lines so that $F$ maps each small square in $W_1$ or in $W_2$.
Then we can build the lifting $\tilde F$ square by square, from button up in the order shown in the picture.

When we extending $\tilde F$ over a small square $\square$ the value of $\tilde F$ is already defined at its lower left corner, say $v$.
So we can apply \ref{lem:covering}.
The part of $\square$ on which $\tilde F$ was already defined is one or two sides adjacent to $v$.
In both cases this set is connected and by \ref{lem:covering} the extension agrees with the already constructed part.
Again, the continuity of $\tilde F$ follows from \ref{ex:u-continuous:closed}.\qeds

\subsection*{Degree}

A loop in $\SS^1$ based at $1$ is a continuous map $f\:[0,1]\to\SS^1$
such that $f(0)=f(1)=1$.
Let $\tilde f\:[0,1]\to\RR$ be the unique lift with $\tilde f(0)=0$.
Since $e^{-1}(1)=\ZZ$, we have $\tilde f(1)\in\ZZ$.
The integer $\tilde f(1)$ will be called \index{degree of loop}\emph{degree} of the loop $f$, it will be denoted by $\deg f$.


\begin{thm}{Theorem}\label{thm:pi1S1Z}
The map $[f]\mapsto \deg f$ defines an isomorphism $\pi_1(\SS^1,1)\to\ZZ$.
\end{thm}

\parit{Proof.}
First let us show that $\deg f$ depends only on the
homotopy class of $f$.
In other words the map $[f]\mapsto \deg f$ is well defined.

Suppose $f_0$ and $f_1$ are two homotopic loops in $\SS^1$.
...

It remains to show that $\deg f_0*f_1=\deg f_0+\deg f_1$ for any two loops $f_0$ and $f_1$.
Let $\tilde f_0$ and $\tilde f_1$ be liftings of $f_0$ and $f_1$ such that $\tilde f_0(0)=\tilde f_1(0)=0$.
Recall that $\deg f_0=\tilde f_0(1)$ and $\deg f_1=\tilde f_1(1)$.
Observe that $t\mapsto \deg f_0+\tilde f_1$ is a lifting of $f_1$ s

We proved that degree defines a homomorphism $\pi_1(\SS^1,1)\cong\ZZ$;
it remains to show that this homomorphism is injective and surjective.
Surjectivity follows ???.
\qeds

\section{Fundamental theorem of algebra}

\begin{thm}{Theorem}\label{cor:FTA}
Every non-constant complex polynomial has a complex root.
\end{thm}

\parit{Proof.}
If a polynomial $p$ has no roots, then for each $r\ge 0$ the loop
$t\mapsto p(r\exp(2\pi i t))/|p(r\exp(2\pi i t))|$ lies in $\SS^1$.
For large $r$ this loop is homotopic to $t\mapsto \exp(2\pi i k t)$, where $k$ is
the degree of $p$, while for $r=0$ it is constant; this forces a change of degree,
which is impossible under a homotopy rel endpoints.

